\section{n3. Full normalization compatibility with finite-covered normal sources}
\textit{Supports: ENTIRE qualified G331; own n1; own n2.}

Here R is any positive-dimensional excellent REDUCED local ring, not
necessarily a domain or equidimensional. Let S be its finite normalization,
keep every marked normal local centre, and let
$\mathfrak c=(R:S)\subset R$. G331 gives a regular conductor element,
the full compact quotient $\rho:Y\to X$, singleton fibres on
$X^\circ=\{v:v(\mathfrak c)<\infty\}$, and precisely conductor-quotient
fibres at infinite value. All nonminimal supports are kept. Let
$Y^\circ=\rho^{-1}X^\circ$. The underlying domain-localization argument
and fibre-gap argument of supplied 429 are reconstructed below, not
credited as new results. The extension to reduced targets and normal
normalization centres with the finite-cover continuity of n2 is the
new scope in this note.

We explicitly use a result of the ENTIRE supplied grant, not a fresh
density source application: G331 n6 proves that support-zero divisorial
valuations are dense on the full section of every positive-dimensional
excellent local domain. Its Rees-DVR/birational argument and the precise
published dependencies remain qualified as in its ledger. A regular
conductor element is nonzero on each normalization domain centre and has
finite value at every such divisorial point. Therefore $Y^\circ$ is dense
in Y, by applying that GRANTED theorem on each clopen centre. This is
not a density inference merely from the word normal, nor a new density
proof or new source-reading credit. The same conclusion in the smooth
closed domain also follows from the supplied translated-Gauss argument,
which is retained as DATA but is not needed as an additional premise.

For finite continuous signed $f$ on X and $F=\rho^*f$, the canonical
lemma reconstructed in n2 applies without target normality: it uses only
G331 comparison, compactness, rational primary density and grading. Thus
both raw C/D suprema on each ring equal the canonical suprema.
Pulling target competitors up always gives
\[
                 \rho^*p_R(f)\le p_S(F)\quad\text{on Y}.    \tag{12}
\]
On EVERY finite-conductor point there is equality, without source
continuity. Indeed let $J=A_m^S(F)$, $q=-w(J)/m\le F$ and choose an
integer $M\ge\max(1,\|q\|_\infty,\max(0,-\min f))$.
For every integer $r\ge1$ the actual downstairs ideal
\[
              B_r=\mathfrak c J^r+\mathfrak a^{mrM}\subset R
\]
is finitely generated, unit-or-primary. The containment uses that the
conductor is an S-ideal contained in R, including every component of S.
Its extended S-ideal is exactly the displayed sum. Consequently, using
only finite ideal generators and their nonnegative coefficients,
\[
 -w(B_r)/(mr)
       =\max\{q(w)-w(\mathfrak c)/(mr),-M\}\le F(w).        \tag{13}
\]
This includes $w(\mathfrak c)=\infty$, where only the truncated term
survives, and involves no subtraction of two infinities. Surjectivity
makes it a target minorant. The proper case has actual negative effective
ample Cartier presentation on its ACTUAL primary blowup and positive
degree on EVERY contracted integral curve. At a finite-conductor point
(13) tends to $q(w)$. Supremizing over m proves equality in (12) on
$Y^\circ$. No uniform conductor error, generic-support deletion or psh
status for a raw increasing supremum is asserted.

Suppose now that $H=p_S(F)$ is continuous. Then the exact criterion is
\[
 \begin{gathered}
 \text{a greatest target C or D minorant exists}\quad\Longleftrightarrow
 \quad H\text{ is constant on EVERY full fibre of }\rho.     \tag{14}
 \end{gathered}
\]
If fibre-constant, H is in $C_S$ by the canonical continuity criterion.
G331 full-fibre POSITIVE descent gives $h\in C_R$ with $\rho^*h=H$.
It is below f and dominates every target competitor by (12).
Conversely a greatest target minorant is continuous by n2's canonical
criterion. Its pullback equals H on the dense $Y^\circ$ by (13), hence
everywhere by continuity. That proves full-fibre constancy. The quotient
identification and positive descent used here are the full G331 results,
including infinite-conductor supports, not just topological descent.

If two lifts of v have $H(w_+)-H(w_-)=\Delta>0$, choose a net
$w_i\in Y^\circ$ tending to $w_+$. Its restrictions $v_i$ tend to v,
and (12)--(13) imply
\[
 p_R(f)(v)\le H(w_-),\qquad p_R(f)(v_i)\longrightarrow H(w_+).
\]
Thus the raw supremum has a non-USC jump at least $\Delta$.
Every actual D minorant g is USC and satisfies
$\limsup_i g(v_i)\le g(v)\le H(w_-)$, also when g(v) is bottom.
Therefore
\[
       \sup_X(p_R(f)-g)\ge\Delta,
       \qquad \|p_R(f)-h_m^R\|_\infty\ge\Delta
       \quad\text{for EVERY degree }m.                     \tag{15}
\]
This also excludes finite generation of the canonical algebra: finitely
many homogeneous generators of degrees $d_j$ give a profile
$-\min_j v(b_j)/d_j=-v((b_j^{D/d_j}))/D$ for common multiple D.
Their degree-one generators generate $A_1$, which contains a boundary
power, so that displayed ideal is unit-or-primary. Its continuous profile
would be the raw supremum, contradicting (15). Failure of finite
generation by itself has not been used to infer discontinuity.
If source continuity is known for all pulled obstacles, a uniform
perturbation of size $\epsilon<\Delta/2$ leaves a gap at least
$\Delta-2\epsilon$, by the Lipschitz bound for the source raw supremum.

The NEW finite-cover application of (14) is the following. For every
normalization centre $S_i$, suppose there is an actual finite domain
extension $S_i\hookrightarrow T_i$, retaining all its closed centres,
such that $T_i$ has continuous geometric envelopes for finite continuous
obstacles. Own n2 then proves source continuity on $S_i$ without a group
action or flatness. The finite product of these local envelopes is H on S:
G331's semilocal gluing identifies C and D with their centre products.
Thus (14) holds for EVERY continuous signed obstacle on the reduced R,
with arbitrary full conductor geometry and these covered normal centres.
The entire criterion is obstacle-specific. We do not prove that every
nonsingleton normalization fibre supplies an obstacle with a discrepancy,
or give an intrinsic necessary-and-sufficient ring classification.

Primary changes divide H and every fibre gap by
$d(v)=v(\mathfrak b)>0$, which is constant on the full fibre;
the gap at v becomes $\Delta/d(v)$. Completion uses the exact G331
boundary and primary-ideal transport, including reductions of any
completed nilpotents. Normal R has unit conductor and singleton
normalization fibres; this criterion then reduces to its source
continuity hypothesis. Zero-dimensional reductions are empty and direct.
The min(f,0) reduced-curve comparison and literal radial distinction are
recorded in n4, not replaced by this conductor criterion.

Artifacts: notes/n3.tex, P12.3-normal-finite-cover-transfer.tex,
source-application-reading-ledger.md, scoped-verdict.md and the ENTIRE
retained 429 compatibility proof. Rerun: python assemble\_delivery.py;
python verify\_delivery.py. Previous 429 localization and gap conclusions
are reconstructed foundations; the new scope is finite-covered normal
sources and reduced targets. No full-original closure is claimed.
